%- block legend
This is an interactive problem.

The jury has secretly chosen a hidden integer $x$ such that $1 \le x \le n$.
Your task is to determine the value of $x$ by asking at most $25$ queries.

In a single query, you pick an integer $q$ ($1 \le q \le n$) and ask the jury how
$x$ compares to it.
%- endblock

%- block interaction
First, read a single integer $n$ ($\VAR{N.min} \le n \le \VAR{N.max | sci}$),
the size of the search range.

Then, you may make at most $25$ queries. To make a query, print a single
integer $q$ ($1 \le q \le n$) on its own line. After each query, read the jury's
response, which is one of the following two strings:

\begin{itemize}
	\item \texttt{<} --- meaning $x < q$;
	\item \texttt{>=} --- meaning $x \ge q$.
\end{itemize}

When you are sure about the value of $x$, print \texttt{! x} on its own line and
terminate your program. The answer line does not count towards the query limit.

After printing each query (or the final answer), do not forget to print a newline
and flush the output buffer. Otherwise, you may receive a verdict of
\textit{Idleness limit exceeded}.
%- endblock

%- block notes
In the first sample, the hidden number is $x = 2$ within the range $[1, 10]$.
The solution narrows the range down with three queries before answering.

In the last sample, the range contains a single value, so the answer can be
given immediately without any query.
%- endblock

%- block editorial
Maintain a candidate interval $[lo, hi]$, initially $[1, n]$. While $lo < hi$,
query the midpoint $q = \lfloor (lo + hi + 1) / 2 \rfloor$. If the response is
\texttt{<}, then $x < q$, so set $hi = q - 1$; otherwise $x \ge q$, so set
$lo = q$. After at most $\lceil \log_2 n \rceil \le 25$ queries the interval
collapses to a single value, which is the answer.
%- endblock
